\begin{abstract}
Transient execution side-channel attacks, such as Spectre, have been shown to break almost all isolation primitives.
We introduce a new security property we call \emph{relaxed microarchitectural isolation} (RMI) that allows \emph{sensitive} programs that are \emph{not-constant-time} to \emph{share memory} with an attacker while restricting the information leakage to that of non-speculative execution.
Although this type of speculative security property is typically challenging to enforce, we show that we can leverage the \emph{enclave} setup to achieve it.
In particular, we use microarchitectural isolation to restrict attacker's observations in conjunction with straightforward hardware mechanisms to limit speculation.
This new design point presents a compelling trade-off between security, usability, and performance, making it possible to efficiently enforce RMI for \emph{any} program.

We demonstrate our approach by implementing and evaluating two simple defense mechanisms that satisfy RMI: (1) \emph{Safe mode}, which disables speculative accesses to shared memory, and (2) \emph{Burst mode}, a localized performance optimization that requires simple program analysis on small code snippets.
Our end-to-end prototype, Citadel, consists of an FPGA-based multicore processor that boots Linux and runs secure applications, including cryptographic libraries and private inference, with less than 5\% performance overhead.
\end{abstract}